\section*{Bab \ref{ch:g7:states-of-matter} --- Wujud Zat dan Perubahan Wujud}

\begin{solution}{exo:g7:states-of-matter:1}
Padat: \omterm{def:g7:states-of-matter:molecule}{molekulnya} terkemas dalam tumpukan yang tertib, bergetar di
tempatnya. Cair: \omterm{def:g7:states-of-matter:molecule}{molekulnya} bersentuhan tetapi tidak tertib, meluncur
melewati satu sama lain. \omterm{def:g3:solids-liquids-gases:gas}{Gas}: \omterm{def:g7:states-of-matter:molecule}{molekulnya} berjauhan, terbang cepat ke
segala arah.
\end{solution}

\begin{solution}{exo:g7:states-of-matter:2}
Pada \omterm{def:g3:solids-liquids-gases:solid}{zat padat} setiap \omterm{def:g7:states-of-matter:molecule}{molekul} menjaga posnya di dalam tumpukannya,
sehingga keseluruhannya mempertahankan bentuknya. Pada \omterm{def:g3:solids-liquids-gases:liquid}{zat cair}
\omterm{def:g7:states-of-matter:molecule}{molekulnya} meluncur dengan bebas melewati satu sama lain sambil tetap
bersentuhan: kerumunannya merosot menyesuaikan wadahnya, sampai tidak
ada yang dapat meluncur lebih rendah --- itulah permukaan yang
mendatar.
\end{solution}

\begin{solution}{exo:g7:states-of-matter:3}
\omterm{def:g2:air-around-us:air}{Udara} sebagian besar berupa ruang kosong di antara \omterm{def:g7:states-of-matter:molecule}{molekul} yang
beterbangan: mendorong mengemas para penerbangnya lebih rapat (dan
hantaman mereka melawan --- itulah pegasnya). \omterm{def:g7:states-of-matter:molecule}{Molekul} air sudah
bersentuhan; tidak ada kekosongan yang dapat diperas keluar, sehingga
penghisapnya bertemu dinding.
\end{solution}

\begin{solution}{exo:g7:states-of-matter:4}
Bahwa \omterm{def:g7:states-of-matter:molecule}{molekul} \omterm{def:g3:solids-liquids-gases:gas}{gas} terbang berjauhan dengan kekosongan yang besar di
antaranya --- beberapa ruangan penuh dapat dikemas menjadi satu. Air,
yang \omterm{def:g7:states-of-matter:molecule}{molekulnya} sudah bersentuhan, tidak dapat \omterm{prop:g7:states-of-matter:squeeze}{dimampatkan}: tidak ada
pompa yang mengemasnya lebih rapat.
\end{solution}

\begin{solution}{exo:g7:states-of-matter:5}
Kehangatan membuat \omterm{def:g7:states-of-matter:molecule}{molekul} yang bertumpuk itu bergetar makin keras dan
makin keras; pada \omterm{def:g4:changes-of-state:melting-point}{titik leburnya} \omterm{def:g3:sound-around-us:vibration}{getarannya} mengalahkan cengkeraman
tumpukannya, barisannya patah, dan \omterm{def:g7:states-of-matter:molecule}{molekulnya} mulai meluncur: \omterm{def:g3:solids-liquids-gases:solid}{zat
padatnya} telah menjadi \omterm{def:g3:solids-liquids-gases:liquid}{zat cair}.
\end{solution}

\begin{solution}{exo:g7:states-of-matter:6}
Perubahan wujud hanya menyusun ulang \omterm{def:g7:states-of-matter:molecule}{molekul} yang sama --- tidak ada
yang diciptakan atau dimusnahkan, sehingga \omterm{def:g3:measuring-mass:mass}{massanya} tidak mungkin
berubah. Kehilangan \omterm{def:g3:measuring-mass:mass}{massa} akan menuntut \omterm{def:g7:states-of-matter:molecule}{molekulnya} lenyap (atau lolos
dari wadahnya --- tanda bahaya palsu milik botol yang tidak tertutup).
\end{solution}

\begin{solution}{exo:g7:states-of-matter:7}
Kelincahan gerak \omterm{def:g7:states-of-matter:molecule}{molekul}. Ketika yang panas menyentuh yang dingin,
\omterm{def:g7:states-of-matter:molecule}{molekul} yang lincah menghantam \omterm{def:g7:states-of-matter:molecule}{molekul} yang lamban, sambil menyerahkan
kelincahannya tumbukan demi tumbukan --- yang lincah menjadi tenang,
yang lamban menjadi cepat --- sampai kedua kerumunannya berbagi satu
langkah: \omterm{def:g3:temperature-thermometers:temperature}{suhu} yang sama.
\end{solution}

\begin{solution}{exo:g7:states-of-matter:8}
\omterm{def:g7:states-of-matter:molecule}{Molekul} parfumnya terbang, bertumbukan dengan \omterm{def:g7:states-of-matter:molecule}{molekul} \omterm{def:g2:air-around-us:air}{udara},
memantul, lalu mengembara --- zigzag demi zigzag menyeberangi
ruangannya. Embusan \omterm{def:g2:air-around-us:wind}{angin} membawa seluruh kerumunan mereka
sekaligus: menunggangi \omterm{def:g2:air-around-us:wind}{angin} mengalahkan terhuyung-huyung menembus
kerumunannya.
\end{solution}

\begin{solution}{exo:g7:states-of-matter:9}
Hanya \omterm{def:g7:states-of-matter:molecule}{molekul} yang \emph{tercepat} yang dapat merobek diri dari
cengkeraman \omterm{def:g3:solids-liquids-gases:liquid}{zat cairnya}, sehingga \omterm{def:g6:water-states:changes}{penguapan} adalah pajak atas
kelincahan: para pelariannya membawa pergi lebih daripada jatahnya,
dan kerumunan yang tertinggal, secara rata-rata, lebih lambat ---
artinya lebih sejuk. Keringat dan jari yang basah menyejuk dengan
mengekspor yang terlincah.
\end{solution}

\begin{solution}{exo:g7:states-of-matter:10}
Di bawah matahari \omterm{def:g7:states-of-matter:molecule}{molekul} yang terperangkap itu terbang lebih cepat
lalu menghantam kulitnya lebih keras dan lebih sering: badainya
menggembungkan balonnya sampai kulit yang teregang mengimbanginya. Di
tempat yang dingin \omterm{def:g7:states-of-matter:molecule}{molekul} yang sama melambat, hantamannya melemah,
dan dorongan \omterm{def:g2:air-around-us:air}{udara} di luarnya menghimpit balonnya menjadi lebih kecil.
\end{solution}

\begin{solution}{exo:g7:states-of-matter:11}
Air yang \omterm{def:g2:ice-water-steam:meltfreeze}{membeku} mengunci diri ke dalam tumpukan yang luar biasa
lapang dan terbuka --- ruang kosongnya lebih banyak daripada yang tadi
dimiliki kerumunan yang meluncur --- sehingga es memakai ruang lebih
banyak untuk \omterm{def:g7:states-of-matter:molecule}{molekul} yang sama: \omterm{def:g3:measuring-mass:mass}{massa} jenisnya di bawah \omterm{def:g6:density-floating:density}{massa jenis}
\omterm{def:g3:solids-liquids-gases:liquid}{zat cairnya}, dan ia \omterm{def:g1:floating-and-sinking:words}{terapung}. Seandainya air biasa saja, es akan
menumpuk rapat, \omterm{def:g1:floating-and-sinking:words}{tenggelam} selagi terbentuk, dan kolam akan \omterm{def:g2:ice-water-steam:meltfreeze}{membeku}
dari dasar ke atas --- tanpa sisa air yang terlindung bagi ikannya.
\end{solution}

\begin{solution}{exo:g7:states-of-matter:12}
Tidak ada yang mengaduk airnya, namun tintanya bepergian: \omterm{def:g7:states-of-matter:molecule}{molekul} \omterm{def:g3:solids-liquids-gases:liquid}{zat
cairnya} sendiri bergerak tanpa henti, mendorong-dorong \omterm{def:g7:states-of-matter:molecule}{molekul}
tintanya langkah acak demi langkah acak menembus kerumunannya ---
kedua \omterm{def:g3:solids-liquids-gases:liquid}{zat cairnya} hidup oleh gerak pada setiap saat, bahkan ketika
gelasnya tampak sempurna tenang. Kehangatan mempercepat tiap
dorongannya, sehingga penyebarannya bergegas. Tinta yang sabar dengan
demikian membuat tarian yang tak terlihat menjadi terlihat --- bukti
yang meyakinkan para peragu tentang \omterm{def:g7:states-of-matter:molecule}{molekul} itu sendiri.
\end{solution}

\begin{solution}{pb:g7:states-of-matter:1}
\textbf{1.} Dalam barisan yang rapi, bahu bertemu bahu, masing-masing
bergoyang di tempatnya --- dan tidak seorang pun boleh bertukar tempat
atau meninggalkan posnya.
\textbf{2.} Barisannya larut: para muridnya tetap bahu bertemu bahu
(bersentuhan!) tetapi kini menganyam dan meluncur melewati satu sama
lain. Mereka masih harus menjaga sentuhan dengan kerumunannya ---
tidak seorang pun boleh berlari bebas.
\textbf{3.} Para muridnya berhamburan, berlari cepat dan lurus sampai
mereka bertumbukan atau menabrak dinding, lalu memantul tanpa
henti --- dan karena tidak ada yang menahan mereka bersama, hanya
dindingnya yang menghentikan mereka: kerumunannya menuntut setiap
sudut yang ditawarkan gedungnya.
\textbf{4.} Ya --- ketiga puluh yang sama pada setiap babaknya: \omterm{def:g3:measuring-mass:mass}{massa}
selamat melewati perubahan wujud, yang dipentaskan.
\textbf{5.} Babak gasnya mengizinkannya --- para pelarinya berdesak
masuk ke separuh gedungnya (dan menghantam para penggiringnya lebih
keras). Desakan \omterm{def:g3:solids-liquids-gases:liquid}{zat cairnya} tidak dapat menyusut: para muridnya sudah
bersentuhan. Yang dipentaskan: \omterm{def:g3:solids-liquids-gases:gas}{gas} dapat diperas, \omterm{def:g3:solids-liquids-gases:liquid}{zat cair} tidak.
\textbf{6.} Bergoyang makin keras dan makin keras di posnya;
babaknya menjadi \omterm{def:g3:solids-liquids-gases:liquid}{zat cair} pada saat goyangannya membongkar barisannya
dan para muridnya mulai meluncur melewati satu sama lain.
\textbf{7.} \omterm{def:g5:everyday-energy:energy}{Energi} tombolnya dibelanjakan untuk memecahkan cengkeraman
tumpukannya --- barisan demi barisan --- bukan untuk gerak yang lebih
cepat: pengkritik \omterm{def:g3:temperature-thermometers:temperature}{termometernya} harus melaporkan kelincahan rata-rata
yang \emph{tidak berubah} sampai barisan yang terakhir patah. Dataran
\omterm{def:g3:temperature-thermometers:temperature}{suhunya}, yang dipentaskan.
\textbf{8.} Murid yang \emph{tercepat} --- hanya merekalah yang dapat
merobek diri. Kelincahan rata-rata desakan yang tersisa turun: \omterm{def:g3:solids-liquids-gases:liquid}{zat
cairnya} menyejukkan dirinya sendiri dengan \omterm{def:g2:ice-water-steam:evapcond}{menguap}, tepat rasa sejuk
milik keringat dan jari yang basah.
\textbf{9.} Para pemerannya terus bertumbukan dengan talinya lalu
mendorongnya ke luar; pemeran yang lebih cepat mendorong lebih keras
dan lebih sering, dan lingkarannya menggembung --- itulah balon di
bawah matahari.
\textbf{10.} ``Bertumpuklah dalam susunan yang \emph{terbuka} ---
sepanjang lengan, lapang'': \omterm{def:g3:solids-liquids-gases:solid}{zat padat} air memakai ruang lebih banyak
daripada \omterm{def:g3:solids-liquids-gases:liquid}{zat cairnya}, sehingga rakit babak esnya akan \omterm{def:g1:floating-and-sinking:words}{terapung} di atas
kerumunan babak \omterm{def:g3:solids-liquids-gases:liquid}{zat cairnya}.
\textbf{11.} Baju birunya didorong-dorong --- disenggol langkah acak
demi langkah acak menembus kerumunan yang meluncur sampai ia
mengembara ke mana-mana: itulah tinta yang membaur. Naikkan tombolnya
dan tiap dorongannya mempercepat, sehingga pengembaraannya bergegas.
\textbf{12.} Misalnya: ``Sandiwara malam ini mementaskan satu gagasan:
semua materi adalah \omterm{def:g7:states-of-matter:molecule}{molekul} yang bergerak. Kalau itu diterima, dua
hukum lama datang cuma-cuma --- \omterm{def:g3:measuring-mass:mass}{massa} yang selamat melewati setiap
perubahan wujud, dan dataran \omterm{def:g6:water-states:changes}{peleburannya}. Yang tidak dapat kami
pentaskan adalah ukurannya: pemeran yang sesungguhnya lebih kecil
daripada penglihatan mana pun dan lebih banyak daripada seluruh
penonton di dunia.''
\end{solution}
